% SYNTHETIC TEST UNIT (integration tools) -- deliberate errors are marked "DELIBERATE"
\chapter{一般方法（合成测试）}\label{ch:method}

\section{Ward 恒等式}\label{sec:method-ward}

\begin{lemma}[多胶子 Ward 恒等式 \tagA]\label{lem:method-ward-multi}
合成测试。（本引理只涉及部分子态。）
\uses{def:dis-kinematics}
\end{lemma}

\paragraph{地位与缺口}全阶论证；被明确假定的一步是具名缺口\cref{asm:appN-brst-nonlocal}（[C]）。

\begin{theorem}[因子化模板 \tagA]\label{thm:method-template}
合成测试。（本定理只涉及部分子态。）
\uses{lem:method-ward-multi,asm:appN-collins,asm:appN-values}
\end{theorem}
\begin{proofsketch}略。\end{proofsketch}

\paragraph{地位与缺口}全阶论证；瓶颈\cref{asm:appN-collins}；具名缺口\cref{asm:appN-brst-nonlocal}。

% a direct user of lem:method-ward-multi: passes C-P17 through the lemma's own status paragraph (spec wording)
\begin{lemma}[中继引理 \tagA]\label{lem:method-relay}
合成测试。（本引理只涉及部分子态。）
\uses{lem:method-ward-multi}
\end{lemma}

\paragraph{地位与缺口}全阶论证；瓶颈\cref{lem:method-ward-multi}。

% DELIBERATE C-P17 W (closure contains lem:method-ward-multi; neither its paragraph nor that of its premise names the
% named gap) and C-P16 W (the paragraph names none of the bottleneck labels)
\begin{proposition}[未点名缺口的命题 \tagA]\label{prop:method-gapless}
合成测试。（本命题只涉及部分子态。）
\uses{lem:method-relay}
\end{proposition}

\paragraph{地位与缺口}全阶论证。
